\documentclass[11pt]{amsart}
\usepackage{amsmath,amssymb,amsthm}
\usepackage[margin=1.05in]{geometry}

\numberwithin{equation}{section}
\newtheorem{theorem}{Theorem}[section]
\newtheorem{proposition}[theorem]{Proposition}
\newtheorem{lemma}[theorem]{Lemma}
\newtheorem{corollary}[theorem]{Corollary}
\theoremstyle{definition}
\newtheorem{definition}[theorem]{Definition}
\theoremstyle{remark}
\newtheorem{remark}[theorem]{Remark}
\newtheorem{example}[theorem]{Example}

\newcommand{\Z}{\mathbb{Z}}
\newcommand{\Q}{\mathbb{Q}}
\newcommand{\R}{\mathbb{R}}
\newcommand{\N}{\mathbb{N}}
\newcommand{\C}{\mathbb{C}}
\newcommand{\cA}{\mathcal{A}}
\newcommand{\cL}{\mathcal{L}}
\newcommand{\cH}{\mathcal{H}}
\newcommand{\cE}{\mathcal{E}}
\newcommand{\pp}{\mathfrak{p}}
\newcommand{\qq}{\mathfrak{q}}
\newcommand{\aaa}{\mathfrak{a}}
\newcommand{\bb}{\mathfrak{b}}
\newcommand{\Ns}{\mathrm{N}}
\newcommand{\Db}{D_{\mathrm{base}}}
\newcommand{\Dv}{D_{\mathrm{vert}}}
\newcommand{\Df}{D_{\mathrm{full}}}
\newcommand{\nusym}{\nu_{\mathrm{sym}}}
\newcommand{\id}{\mathrm{id}}

\PassOptionsToPackage{hyphens}{url}
\usepackage[hidelinks,bookmarks=false]{hyperref}

\title[Twists, detailed balance, and tensorisation]{Quantum optimal transport on Bost--Connes systems:\\ the twist is harmful, the generator has a sign error,\\ and the transition is invisible to local inequalities}
\author{Ruqing Chen}
\address{GUT Geoservice Inc., Montreal, Canada}
\email{ruqing@hotmail.com}
\thanks{Sources, code and data for the entire series: \url{https://github.com/Ruqing1963/localized-bost-connes}; this paper is in \texttt{papers/XVI-transport/}.}
\date{}

\begin{document}

\begin{abstract}
We examine three proposals for a quantum optimal transport theory on the type $\mathrm{III}$
Bost--Connes systems and find that each fails, in each case for an identifiable and
correctable reason. Two of the corrections are clean; the third is a structural obstruction.

\emph{The modular twist is not merely unnecessary but harmful.} With
$\Dv\epsilon_\aaa=(\log\Ns\aaa)\epsilon_\aaa$ and $\vartheta=\sigma_{-i\beta}$, one computes
\[
[\Dv,\mu_\pp]_\vartheta\,\epsilon_\bb=\bigl(\log\Ns\pp+(1-\Ns\pp^{\beta})\log\Ns\bb\bigr)\epsilon_{\pp\bb},
\]
which is \emph{unbounded}. The untwisted commutator is bounded and elementary:
$\Dv\mu_\pp=\mu_\pp(\Dv+\log\Ns\pp)$ gives $[\Dv,\mu_\pp]=(\log\Ns\pp)\mu_\pp$. The reason is
that a twist repairs a \emph{multiplicative} mismatch, whereas $\Dv$ is \emph{additive} in
$\log\Ns$; introducing $\vartheta$ manufactures the mismatch it was meant to cure. The
untwisted $\Df=\Db\otimes\sigma_1+\Dv\otimes\sigma_2$ does have $\ker L=\C1$ --- and, as it
turns out, exactly when $S$ is norm-separated, the same condition that governs the
equal-norm analysis of Paper XIV of the series (its Theorem 1.3). Whether it is a spectral
triple reduces to whether the horizontal operator has compact resolvent, which is not
established.

\emph{The proposed Lindbladian does not preserve the KMS state.} Writing $a_\pp,b_\pp$ for
the rates of $\mu_\pp^*(\cdot)\mu_\pp$ and $\mu_\pp(\cdot)\mu_\pp^*$, the KMS condition gives
$\varphi(\cL f)=\sum_\pp\{\nu(fe_\pp)(a_\pp\Ns\pp^{\beta}-b_\pp)+\nu(f)(b_\pp\Ns\pp^{-\beta}-a_\pp)\}$,
which vanishes for all $f$ if and only if $b_\pp=a_\pp\Ns\pp^{+\beta}$. The proposal takes
$b_\pp=a_\pp\Ns\pp^{-\beta}$: \emph{the sign of the exponent is wrong}, and with it
$\varphi(\cL e_\pp)=(1-t^2)(1-t)\ne0$ for $t=\Ns\pp^{-\beta}$. The corrected ratio
$b_\pp/a_\pp=e^{\beta\log\Ns\pp}$ is the detailed-balance factor, as it must be.

\emph{Critical slowing down does not occur.} The proposed asymptotic
$\alpha(\beta)\asymp\Pi_\infty(\beta)=\prod_{\ell}\frac{1-\ell^{-\beta}}{1+\ell^{-\beta}}$
is of the wrong shape: spectral gaps, modified log-Sobolev constants and Bakry--\'Emery
curvature bounds \emph{tensorise to an infimum} of single-site quantities, each a function
of its own $\ell^{-\beta}$ alone, not to a product over $\ell$. The single-site defects
$1-\ell^{-\beta}$ are bounded below by $1-3^{-\beta}\ge0.66$, and the single-site spectral
gap by $(1-t)^2/4t$ (the Poincar\'e constant of the geometric law), so with unit rates the
gap stays bounded below near $\beta_c$ while $\Pi_\infty\to0$. The arithmetic transition is
a \emph{tail} phenomenon --- the Kakutani singularity of an infinite product --- and such
tensorising local quantities cannot see it. This is the same mechanism that made the metric
transition of Paper IV of the series (its Theorem 6.2) continuous rather than first-order.
\end{abstract}

\maketitle

\section{The twist}

We use the notation of \cite{Main}: $\cA_{\Q,S}=C(Y_S)\rtimes J_S$, spanned by the monomials
$\mu_\aaa f\mu_\bb^*$, with $\sigma_t(\mu_\aaa)=\Ns\aaa^{it}\mu_\aaa$; $X$ and $\nusym$ are
the compact model and its symmetric measure. Papers III, IV and XIV of the series, referred
to by numeral, are unpublished companions; every fact used from them is restated or
re-derived in place. Represent $\cA_{\Q,S}$ on $\cH_0=L^2(X,\nusym)\otimes\ell^2(J_S)$ in the Toeplitz
representation, $\mu_\aaa\epsilon_\bb=\epsilon_{\aaa\bb}$, and set
$\Dv\epsilon_\aaa=(\log\Ns\aaa)\epsilon_\aaa$. Let $\vartheta=\sigma_{-i\beta}$, so
$\vartheta(\mu_\aaa)=\Ns\aaa^{\beta}\mu_\aaa$ and $\vartheta(f)=f$ for $f\in C(Y_S)$.

\begin{proposition}[The twisted commutator is unbounded]\label{prop:unbounded}
For every prime $\pp$ and every $\beta>0$,
\begin{equation}\label{eq:twist}
[\Dv,\mu_\pp]_\vartheta\,\epsilon_\bb
=\bigl(\Dv\mu_\pp-\vartheta(\mu_\pp)\Dv\bigr)\epsilon_\bb
=\Bigl(\log\Ns\pp+\bigl(1-\Ns\pp^{\beta}\bigr)\log\Ns\bb\Bigr)\epsilon_{\pp\bb},
\end{equation}
whose coefficient tends to $-\infty$ as $\Ns\bb\to\infty$. Hence
$[\Dv,\mu_\pp]_\vartheta\notin\mathcal B(\cH_0)$.
\end{proposition}

\begin{proof}
$\Dv\mu_\pp\epsilon_\bb=\log\Ns(\pp\bb)\epsilon_{\pp\bb}=(\log\Ns\pp+\log\Ns\bb)\epsilon_{\pp\bb}$
and $\vartheta(\mu_\pp)\Dv\epsilon_\bb=\Ns\pp^{\beta}(\log\Ns\bb)\epsilon_{\pp\bb}$;
subtract. Since $\Ns\pp^\beta\ne1$ the coefficient is unbounded in $\bb$.
\end{proof}

\begin{proposition}[The untwisted commutator is bounded]\label{prop:bounded}
$\Dv\mu_\aaa=\mu_\aaa\bigl(\Dv+\log\Ns\aaa\bigr)$, hence
\[
[\Dv,\mu_\aaa]=(\log\Ns\aaa)\,\mu_\aaa,\qquad \bigl\|[\Dv,\mu_\aaa]\bigr\|=\log\Ns\aaa<\infty,
\]
and $[\Dv,f]=0$ for $f\in C(Y_S)$.
\end{proposition}

\begin{proof}
Both sides send $\epsilon_\bb$ to $(\log\Ns\aaa+\log\Ns\bb)\epsilon_{\aaa\bb}$; the vanishing
on $C(Y_S)$ is the observation recorded in Paper IV of the series (its Proposition 2.1),
$\Dv$ being diagonal in a basis with respect to which $C(Y_S)$ acts diagonally.
\end{proof}

\begin{remark}[Why the twist backfires]\label{rem:whybackfire}
A Connes--Moscovici twist \cite{CM} repairs a \emph{multiplicative} discrepancy: it is used when
conjugation by the algebra rescales $D$, so that $Dx-xD$ fails to be bounded by a factor,
and $\vartheta$ absorbs the factor. Here $\Dv$ is \emph{additive} in $\log\Ns$, so
$Dx-xD$ is already bounded, and inserting the multiplicative $\vartheta(\mu_\aaa)=\Ns\aaa^\beta\mu_\aaa$
introduces a factor where there was none. Equation \eqref{eq:twist} is precisely the
manufactured mismatch. The twisted triple, with twist $\sigma_{i\beta}$, was the candidate
proposed in the closing assessment of Paper IV; this section is the reason we do not pursue
it.
\end{remark}

\begin{remark}[Numerical illustration]\label{rem:numtwist}
For $\Ns\pp=3$, $\beta=1$ the coefficient in \eqref{eq:twist} takes the values
$-3.5,\,-12.7,\,-26.5,\,-54.2$ at $\Ns\bb=10,10^3,10^6,10^{12}$; for $\beta=2$,
$-17.3,\,-54.2,\,-109.4,\,-219.9$. The untwisted commutator has constant norm
$\log3=1.0986$.
\end{remark}

\section{The untwisted operator, and the reappearance of norm-separation}

Let $\Df=\Db\otimes\sigma_1+\Dv\otimes\sigma_2$ on $\cH=\cH_0\otimes\C^2$, with $\Db$ the
horizontal operator of Paper IV (its Definition 2.4), amplified to act on the
$L^2(X,\nusym)$ factor of every summand $L^2(X,\nusym)\otimes\epsilon_\bb$, and
$L(x)=\|[\Df,x]\|$. In this representation $\mu_\aaa$ acts on $\ell^2(J_S)$ alone and $f\in
C(Y_S)$ acts on the summand indexed by $\bb$ as multiplication by the translate
$f\circ m_\bb$, so $[\Db,\mu_\aaa]=0$ while $[\Db,f]$ is the $\bb$-indexed family
$[\Db,f\circ m_\bb]$.

\begin{proposition}[Kernel]\label{prop:kernel}
The operators $\sigma_1,\sigma_2$ anticommute and $x$ acts diagonally, so
$[\Df,x]=[\Db,x]\otimes\sigma_1+[\Dv,x]\otimes\sigma_2$, which vanishes iff both terms do. By
Proposition~\ref{prop:bounded}, $\ker[\Dv,\cdot\,]=\cA^\sigma_{\Q,S}$. Hence
\[
\ker L=\bigl\{x\in\cA^\sigma:[\Db,x]=0\bigr\},
\]
which equals $\C1$ if and only if $\cA^\sigma=C(Y_S)$, i.e.\ if and only if $S$ is
norm-separated in the sense of Theorem 1.3 of Paper XIV.
\end{proposition}

\begin{proof}
$[\Dv,\mu_\aaa f\mu_\bb^*]=(\log\Ns\aaa-\log\Ns\bb)\mu_\aaa f\mu_\bb^*$, zero iff
$\Ns\aaa=\Ns\bb$, so the kernel of the vertical commutator is exactly the fixed-point algebra
of $\sigma$. For $f\in C(Y_S)$, $[\Db,f]=0$ forces $[\Db,f\circ m_1]=0$ on the summand
$\bb=(1)$, hence $f\in\C1$ by the Lip-kernel lemma of Paper IV (its Lemma 3.1); if
$\cA^\sigma\supsetneq C(Y_S)$ the element $\mu_\aaa\mu_\bb^*$ with $\Ns\aaa=\Ns\bb$,
$\aaa\ne\bb$, lies in $\ker[\Dv,\cdot]$ and, acting on $\ell^2(J_S)$ alone, also in
$\ker[\Db,\cdot]$.
\end{proof}

\begin{remark}[Spectral triple: not established]\label{rem:notriple}
Since $\Db$ and $1\otimes\Dv$ commute and are placed on anticommuting Clifford directions,
$\Df^2=\Db^2\otimes1+1\otimes\Dv^2$, and on the summand $L^2(X,\nusym)\otimes\epsilon_\aaa$
this is $\Db^2+(\log\Ns\aaa)^2$. Over $\Q$ the norms $\Ns\aaa$ are distinct and tend to
infinity, so $\Df$ has compact resolvent if and only if $\Db$ does on $L^2(X,\nusym)$ ---
a property of the weight sequence of Paper IV, which that paper does not establish and we do
not claim; its own assessment leaves the spectral triple on the crossed product open, and so
do we. In particular no index theory is asserted. The pair $(\cA_{\Q,S},L)$ can still be a
compact quantum metric space in Rieffel's sense \cite{Rieffel} --- which requires only that
the Lipschitz ball be totally bounded modulo scalars --- and it is that notion, not the
spectral triple, that is at issue here; the two are often conflated.
\end{remark}

\begin{remark}[Open Problem 1 of Paper IV]\label{rem:openone}
Proposition~\ref{prop:kernel} settles the kernel question but not the compact quantum metric
space question, which additionally requires total boundedness of
$\{x:L(x)\le1,\varphi(x)=0\}$. The vertical directions have Lipschitz cost bounded below by
$\log\Ns\pp_0$ for the smallest prime, which is encouraging; we have not verified total
boundedness and do not claim it.
\end{remark}

\section{The generator: a sign error in the detailed balance}

\begin{proposition}[Invariance forces the detailed-balance ratio]\label{prop:balance}
Let
\[
\cL(x)=\sum_{\pp\in S}a_\pp\bigl(\mu_\pp^*x\mu_\pp-x\bigr)+\sum_{\pp\in S}b_\pp\Bigl(\mu_\pp x\mu_\pp^*-\tfrac12(e_\pp x+xe_\pp)\Bigr),
\qquad e_\pp=\mu_\pp\mu_\pp^* ,
\]
the GKLS form with jump operators $\mu_\pp$ and $\mu_\pp^*$. Then $\cL(1)=0$, and for
$f\in C(Y_S)$ --- where $\tfrac12(e_\pp f+fe_\pp)=e_\pp f$ --- the KMS condition
$\varphi(xy)=\varphi(y\,\sigma_{i\beta}(x))$ gives
$\varphi(\mu_\pp^*x\mu_\pp)=\Ns\pp^{\beta}\varphi(xe_\pp)$ and
$\varphi(\mu_\pp x\mu_\pp^*)=\Ns\pp^{-\beta}\varphi(x)$ --- the scaling relations recorded
as Proposition 5.1 of Paper III of the series --- whence
\begin{equation}\label{eq:inv}
\varphi\bigl(\cL(f)\bigr)=\sum_{\pp}\Bigl\{\nu(fe_\pp)\bigl(a_\pp\Ns\pp^{\beta}-b_\pp\bigr)
+\nu(f)\bigl(b_\pp\Ns\pp^{-\beta}-a_\pp\bigr)\Bigr\}.
\end{equation}
This vanishes for every $f$ if and only if
\begin{equation}\label{eq:db}
b_\pp=a_\pp\,\Ns\pp^{+\beta}\qquad\text{for every }\pp\in S .
\end{equation}
\end{proposition}

\begin{proof}
$\cL(1)=\sum a_\pp(\mu_\pp^*\mu_\pp-1)+\sum b_\pp(e_\pp-e_\pp)=0$. The two scaling relations
are the KMS condition with $y=\mu_\pp$, respectively $y=\mu_\pp^*$, and
$\sigma_{i\beta}(\mu_\pp)=\Ns\pp^{-\beta}\mu_\pp$. For \eqref{eq:inv} apply them termwise.
If \eqref{eq:db} holds, both brackets vanish. Conversely write
$A_\pp=a_\pp\Ns\pp^\beta-b_\pp$ and $B=\sum_\pp(b_\pp\Ns\pp^{-\beta}-a_\pp)$; under $\nu$
the valuations are independent with $\nu(e_\pp)=t_\pp=\Ns\pp^{-\beta}$ and
$\nu(e_\pp e_{\qq})=t_\pp t_{\qq}$ for $\pp\ne\qq$, so $f=1$ gives
$\sum_\pp A_\pp t_\pp+B=0$ while $f=e_{\qq}$ gives
$t_{\qq}\bigl(A_{\qq}+\sum_{\pp\ne\qq}A_\pp t_\pp+B\bigr)=0$; subtracting,
$A_{\qq}(1-t_{\qq})=0$, so $A_{\qq}=0$ for every $\qq$, and then $B=0$ term by term.
\end{proof}

\begin{corollary}[The proposed generator is not invariant]\label{cor:signerror}
The choice $a_\pp=w_\pp$, $b_\pp=w_\pp\Ns\pp^{-\beta}$ gives
$b_\pp=a_\pp\Ns\pp^{-\beta}$, violating \eqref{eq:db}: the sign of the exponent is
reversed. Explicitly, with $t=\Ns\pp^{-\beta}$ and $w_\pp=1$: for every $\qq\ne\pp$ the
proposed rates satisfy $A_\qq t_\qq+B_\qq=(t_\qq^{-1}-t_\qq)t_\qq+(t_\qq^2-1)=0$, so the
$\qq$-terms of \eqref{eq:inv} cancel at $f=e_\pp$ and the value is carried by the
$\pp$-term alone:
\[
\varphi\bigl(\cL(e_\pp)\bigr)=t\bigl(t^{-1}-t\bigr)+t\bigl(t^2-1\bigr)=(1-t^2)(1-t)>0
\]
for $t\in(0,1)$; numerically $0.5926$ at $(\Ns\pp,\beta)=(3,1)$, $0.8779$ at $(3,2)$,
$0.9033$ at $(5,1.5)$, $0.8397$ at $(7,1)$. With $b_\pp=a_\pp\Ns\pp^{+\beta}$ the same
expression is exactly $0$.
\end{corollary}

\begin{remark}[What the corrected ratio means]\label{rem:dbmeaning}
Condition \eqref{eq:db} reads $b_\pp/a_\pp=e^{\beta\log\Ns\pp}$: the backward rate exceeds
the forward rate by the Boltzmann factor of the energy $\log\Ns\pp$ transferred by
$\mu_\pp$. This is quantum detailed balance in the sense standard for quantum Markov
semigroups \cite{CarlenMaas}, and it is the only choice compatible with
$\varphi_\beta$. The proposal's $\Ns\pp^{-\beta}$ is the anti-detailed-balance choice, which
drives the system away from equilibrium.

Note the price: \eqref{eq:db} requires $\sum_\pp a_\pp\Ns\pp^{\beta}<\infty$ for the second
sum to converge in norm, so the weights must decay faster than $\Ns\pp^{-\beta}$. That is a
genuine constraint but is satisfiable, for instance by $a_\pp=\Ns\pp^{-\beta-1-\epsilon}$.
We have verified invariance, $\varphi\circ\cL=0$; full KMS-symmetry of the semigroup with
respect to $\varphi_\beta$ in the sense of Cipriani \cite{Cipriani} is a stronger statement,
which we have not checked.
\end{remark}

\section{No critical slowing down}

\begin{theorem}[Tensorisation to an infimum]\label{thm:tensorise}
Let $\cL=\sum_\pp a_\pp\cL_\pp$ act on the product system $(V,\pi_\beta)=\bigotimes_\pp(\N,\mu_{t_\pp})$
with $t_\pp=\Ns\pp^{-\beta}$, each $\cL_\pp$ the birth--death generator on its own factor
with the corrected rates $(a_\pp,\,a_\pp t_\pp^{-1})$ of \eqref{eq:db}. Then:
\begin{enumerate}
\item the spectral gap satisfies $\operatorname{gap}(\cL)=\inf_{\pp}a_\pp\operatorname{gap}(\cL_\pp/a_\pp)$, and by the
Poincar\'e inequality for the geometric law (Paper IV, its Lemma 6.1, with sharp constant
$C(t)=4t/(1-t)^2$)
\[
\operatorname{gap}(\cL_\pp/a_\pp)\ \ge\ \frac{(1-t_\pp)^2}{4t_\pp}\ \ge\ \frac{(1-t_\pp)^2}{4};
\]
\item the modified log-Sobolev constant satisfies
$\alpha(\beta)=\inf_{\pp}a_\pp\,\alpha_0(t_\pp)$, where $\alpha_0(t)$ is the single-site
constant of the normalized chain: entropy is subadditive over the product measure and the
Dirichlet form is additive, giving $\ge$, and single-site test functions give $\le$.
\end{enumerate}
Both constants are therefore \emph{infima of single-site quantities, each a function of its
own $t_\pp$ alone} --- not products over $\pp$. In particular: with rates bounded below,
$\inf_\pp a_\pp>0$, the gap stays $\ge\inf_\pp a_\pp\,(1-3^{-\beta})^2/4>0$ uniformly for
$\beta$ in a neighbourhood of $\beta_c$; with the norm-summable rates of
Remark~\ref{rem:dbmeaning} the infima vanish for every $\beta$, for the weight reason
alone and $\beta$-independently. In neither reading does any of these constants acquire the
product shape $\Pi_\infty(\beta)$, whose vanishing at $\beta_c$ is driven by the divergence
of $\sum_\ell\ell^{-\beta}$, a tail quantity invisible to an infimum.
\end{theorem}

\begin{proof}
For commuting single-site generators on a product measure, the Dirichlet form and the
variance resp.\ entropy split: $\operatorname{Var}=\sum$ over sites of averaged single-site
variances bounds the gap from below by the infimum, and single-site eigenfunctions attain
it; the same subadditivity of entropy (a product measure satisfies
$\operatorname{Ent}(f)\le\sum_\pp\mathbb E\operatorname{Ent}_\pp(f)$) with the additive
Dirichlet form gives $\alpha\ge\inf a_\pp\alpha_0(t_\pp)$, and functions of a single
coordinate give $\le$. The single-site gap bound is Lemma 6.1 of Paper IV, whose Dirichlet
form $\mathbb E_{\mu_t}|\Delta h|^2$ is that of the normalized chain with rates $(1,t^{-1})$.
\end{proof}

\begin{corollary}[The proposed asymptotic is of the wrong shape]\label{cor:wrongshape}
$\Pi_\infty(\beta)=\prod_{\ell\in R}\frac{1-\ell^{-\beta}}{1+\ell^{-\beta}}$ is a product over
all $\ell$ and tends to $0$ as $\beta\downarrow\beta_c$, precisely because the series
$\sum_\ell\ell^{-\beta}$ diverges there. By Theorem~\ref{thm:tensorise} the gap and MLSI
constants are infima of single-site quantities, which do not degenerate this way: the
single-site defect $1-\ell^{-\beta}$ never falls below $1-3^{-\beta}\ge0.66$, so with unit
rates the gap infimum stays $\ge(1-3^{-\beta})^2/4\ge0.11$. Numerically, over the primes
$3\le\ell\le10^4$,
\[
\begin{array}{lcc}
\beta & \inf_\ell(1-\ell^{-\beta}) & \Pi_\infty\\\hline
1.5 & 0.8076 & 3.70\cdot10^{-1}\\
1.2 & 0.7324 & 1.27\cdot10^{-1}\\
1.05 & 0.6845 & 3.58\cdot10^{-2}\\
1.001 & 0.6670 & 1.86\cdot10^{-2}
\end{array}
\]
The product collapses; the single-site defect does not fall below $1-3^{-\beta}\ge0.66$.
Hence $\alpha(\beta)\not\asymp\Pi_\infty(\beta)$, and there is no critical slowing down of
this kind.
\end{corollary}

\begin{corollary}[Bakry--\'Emery]\label{cor:BE}
The same argument applies to the curvature-dimension condition: for product systems with
additive generators the tensorised curvature bound is again an infimum of single-site
quantities, each a function of its own $t_\pp$ alone. Whatever the single-site value, the
bound therefore has no product shape, and no degeneration tied to $\beta_c$ beyond what the
weights impose $\beta$-independently: there is no Ricci-flat singularity at the transition.
\end{corollary}

\begin{remark}[Why the transition is invisible]\label{rem:invisible}
The arithmetic phase transition is the Kakutani dichotomy: an infinite product of
single-site affinities degenerates although every factor stays bounded away from $0$
(Paper IV, its Lemma 5.1). It is therefore a \emph{tail} event, detectable only by quantities
sensitive to infinitely many coordinates at once. Modified log-Sobolev constants and
Bakry--\'Emery curvature are local: they tensorise, so they see only the worst single site.
The same mechanism produced Theorem 6.2 of Paper IV, where the metric transition was
continuous rather than first-order because the order parameter is a tail observable outside
the Lipschitz domain, and Theorem 2.3 of Paper III, where $\Xi_\beta$ was identified as an
invariant of a measure class. Three different local invariants, one tail phenomenon, three
failures to detect it.
\end{remark}

\section{Assessment}

\[
\begin{array}{lll}
\text{Proposal} & \text{Status} & \text{Reference}\\\hline
[\Df,x]_\vartheta\ \text{bounded} & \textbf{false: unbounded} & \text{Prop.~1.1}\\
\text{untwisted}\ [\Dv,\mu_\aaa]\ \text{bounded} & \text{true} & \text{Prop.~1.2}\\
\text{the twist is needed} & \textbf{no; it creates the problem} & \text{Rem.~1.3}\\
\ker L=\C1 & \text{true iff }S\ \text{norm-separated} & \text{Prop.~2.1}\\
(\cA,\cH,\Df)\ \text{a spectral triple} & \text{not established: reduces to }\Db & \text{Rem.~2.2}\\
\cA_{\Q,S}\ \text{a compact quantum metric space} & \text{not established} & \text{Rem.~2.3}\\
\varphi\circ\cL_\beta=0\ \text{as proposed} & \textbf{false: sign error} & \text{Cor.~3.2}\\
b_\pp=a_\pp\Ns\pp^{+\beta} & \textbf{the correction} & \text{Prop.~3.1}\\
\text{KMS-symmetry (Cipriani sense)} & \text{not checked} & \text{Rem.~3.3}\\
\alpha(\beta)\asymp\Pi_\infty(\beta) & \textbf{false: infimum, not product} & \text{Thm.~4.1, Cor.~4.2}\\
\text{critical slowing down} & \textbf{does not occur} & \text{Cor.~4.2}\\
\kappa(\beta)\to0\ \text{at}\ \beta_c & \textbf{false} & \text{Cor.~4.3}
\end{array}
\]

Two of the three failures are repairable and the repairs are one-line: drop the twist, and
flip the sign in the detailed-balance ratio. The third is not a repair but a structural
statement, and it is the one worth keeping. Modified log-Sobolev and Bakry--\'Emery
tensorise, hence are determined by the worst single coordinate, hence cannot detect a
Kakutani dichotomy, which is by construction a statement about infinitely many coordinates
whose individual contributions are all harmless; and the Connes distance of Paper IV is
blind to the same transition in the complementary sense of its Theorem 4.2, the tail order
parameter and all its truncations lying outside every Lipschitz domain. A quantum optimal
transport theory that detects the arithmetic transition would have to be built from a
genuinely non-tensorising quantity, and we do not know one.

Three questions. Is $(\cA_{\Q,S},L)$ with the untwisted $\Df$ a compact quantum metric
space, i.e.\ is the Lipschitz ball totally bounded? Does the corrected generator of
Proposition~\ref{prop:balance} satisfy KMS-symmetry in the sense of Cipriani
\cite{Cipriani}, and if so what
is the associated non-tracial carré du champ? And is there any transport-theoretic quantity
that fails to tensorise --- the natural candidates being those built from the tail
$\sigma$-algebra, where by \cite[Prop.~7.2]{SeriesI} the order parameter lives, rather than
from a sum of single-site Dirichlet forms.

\begin{thebibliography}{9}
\bibitem{CarlenMaas} E. A. Carlen, J. Maas, \emph{Gradient flow and entropy inequalities for quantum Markov semigroups with detailed balance}, J. Funct. Anal. 273 (2017), 1810--1869.
\bibitem{Cipriani} F. Cipriani, \emph{Dirichlet forms and Markovian semigroups on standard forms of von Neumann algebras}, J. Funct. Anal. 147 (1997), 259--300.
\bibitem{CM} A. Connes, H. Moscovici, \emph{Type $\mathrm{III}$ and spectral triples}, in: Traces in Number Theory, Geometry and Quantum Fields, Vieweg, 2008, 57--71.
\bibitem{Rieffel} M. A. Rieffel, \emph{Metrics on state spaces}, Doc. Math. 4 (1999), 559--600.
\bibitem{Main} R. Chen, \emph{KMS state uniqueness and phase transitions in localized Bost--Connes systems}, preprint, 2026, \newline\url{https://doi.org/10.5281/zenodo.22152101}.
\bibitem{SeriesI} R. Chen, \emph{Localized Bost--Connes systems I: factor type and the spectrum of transitions}, preprint, 2026, DOI \href{https://doi.org/10.5281/zenodo.22160827}{10.5281/zenodo.22160827}.
\end{thebibliography}

\end{document}
